% !TeX spellcheck = en_US
\documentclass[11pt,parskip=full+]{scrartcl}

\usepackage[utf8]{inputenc}
\usepackage[T1]{fontenc}

\usepackage[a4paper,margin=2cm]{geometry}
\usepackage[american]{babel}
\usepackage[autostyle]{csquotes}
\usepackage{setspace}
\usepackage{enumitem}
\usepackage{calc}
\usepackage[
backend=biber,
style=numeric,
sorting=nyt, % Sort references by name, year, and title.
natbib=true, % Enable natbib compatibility.
sortcites,
%    dashed=false, % Always display author names in the bibliography.
%    uniquelist=false, % Display only the first author (e.g., Author et al., Year).
urldate=iso, % Use ISO format for accessed dates in online sources.
%    seconds=true % Include seconds in date-time fields.
	maxcitenames=1,
]{biblatex}
\usepackage{hyperref}
\hypersetup{
	bookmarksnumbered, % Enables the chapter and section numbers within the PDF bookmarks.
	colorlinks=true, % Enable colored text for links.
	urlcolor=black, % Color for \url and \href links.
	linkcolor=black, % Color for internal links (e.g., \nameref).
	citecolor=black, % Color for citation links.
}
\addbibresource{bibliography.bib}
\usepackage{tabularx}
\usepackage{array}


\title{Global Optimization Methods for Auto-Tuning HPC Particle Simulations}
\subtitle{Exposé for the Interdisciplinary Project}
\author{Niklas Ladurner}
\date{October 20, 2026}

\makeatletter
\renewcommand{\maketitle}{
	\begin{center}
		{\LARGE\bfseries \@title}\\[2pt]
		\large \@author{} --- \@subtitle
	\end{center}
	\let\thedate\@date
	\let\theauthor\@author
	\let\thetitle\@title
}
\makeatother


\begin{document}
\maketitle
\pagenumbering{gobble}
\onehalfspacing


\section*{Project Description}
\vspace*{-1em}
The simulation of particles and their interactions has become an important area of research, with a large and diverse number of applications ranging from materials sciences to drug development. Due to the large number of interacting bodies, such systems can no longer be solved by analytic means, therefore numerical approaches are required. The development and improvement of these methods has thus become a key research direction in Scientific Computing. Accordingly, there exists a number of different particle simulation frameworks, typically optimized for specific applications and therefore with corresponding strengths and weaknesses.

One such particle simulation framework, developed at TUM, is AutoPas. The main feature that sets AutoPas apart from alternative simulation tools is its core concept of Auto-Tuning. Depending on a given scenario, different algorithms and their associated parameters are optimal in terms of energy consumption or runtime of the simulation. Hence, any fixed choice of such a configuration is optimal only for a select number of scenarios and may perform poorly for others. With Auto-Tuning, the configuration choice is flexible: at a set tuning interval, so-called tuning iterations are performed to sample a number of different configurations. This allows for the collection of several metrics for each tested configuration, based on which a tuning strategy can then select the best performing configuration. Ideally, it is a better fit than the previous configuration and is used for computation of the remaining iterations in the tuning interval, minimizing a set optimization criterion such as energy consumption or runtime. \cite{Newcome2025,Gratl2019,Seckler2021,Gratl2021}

Sampling a given configuration entails computing a whole iteration step of the simulation, which is comparably expensive. Adding new parameters or expanding any parameter's range increases the size of the parameter space significantly, thus testing out all possible configurations (full-search) becomes infeasible. To keep Auto-Tuning viable, the number of candidate configurations to sample must be limited by applying more sophisticated tuning strategies. An approach that has been explored recently uses machine learning algorithms for a better selection of configurations to sample. However, the presented strategy uses a generalized model, and improving it for a specific scenario requires scenario-specific training, which in turn requires collection of measurements for the given scenario to be performed beforehand. Additionally, models trained on such measurement data may be sensitive to a number of variables in the training environment, such as code versions or hardware configurations, which limits portability across different systems. \cite{Metscher2025}

This interdisciplinary project intends to explore a different class of optimization methods to develop new tuning strategies that may overcome such drawbacks. The focus will lie on global optimization algorithms, in particular simulated annealing and genetic algorithms. These approaches are probabilistic heuristics that can both be applied on discrete domains to find global minima of non-convex functions, which in our case describes the process of optimizing for the best-fitting configuration. The mathematical foundations of these models constitute the connection to the application subject mathematics, and the previous knowledge I gathered in my undergraduate studies will be very helpful in understanding and adapting these concepts. Additionally, the area of Scientific Computing, particularly numerical simulations, sits itself at the intersection of computer science, physics and mathematics and is therefore a highly interdisciplinary subject. The accompanying lecture on convex optimization aligns well with this interdisciplinary project, as it introduces the main concepts and methods in global optimization which are applied in this work.


\section*{Project Details}
\vspace*{-1em}
%\begin{description}[leftmargin=!,labelwidth=\widthof{\textbf{Supervisor}}, itemsep=-1em]
%	\item[\textbf{Title}] \thetitle{}
%	\item[\textbf{Lecture}] Convex Optimization (0000004749)
%	\item[\textbf{Supervisor}] Univ.-Prof. Dr. Hans-Joachim Bungartz
%	\item[\textbf{Advisor}] Manish Kumar Mishra, M.Sc. (hons)
%\end{description}

\begin{table}[htpb]
	\noindent
	\begin{tabularx}{\textwidth}{@{}p{3.5cm} X@{}}
		\textbf{\textsf{Title}}      & \thetitle{}                                                                     \\
		\textbf{\textsf{Lecture}}    & Convex Optimization (0000004749)                                                \\
		\textbf{\textsf{Supervisor}} & Univ.-Prof. Dr. Hans-Joachim Bungartz                                           \\
		\textbf{\textsf{Advisor}}    & Manish Kumar Mishra, M.Sc. (hons)                                               \\
	\end{tabularx}
\end{table}

\section*{Project Timeline}
\vspace*{-1em}
%\begin{description}[leftmargin=!,labelwidth=\widthof{\textbf{Algorithmic Design}}, itemsep=-1em]
%	\item[\textbf{Literature Review}] October -- November.
%	\item[\textbf{Algorithmic Design}] December. Familiarization with AutoPas' tuning strategies and codebase, design of the simulated annealing and genetic algorithm tuning strategies.
%	\item[\textbf{Pause}] January -- May. Pausing work due to an exchange semester at UC Berkeley.
%	\item[\textbf{Implementation}] June -- July.
%	\item[\textbf{Documentation}] August. Evaluation of the implemented approaches, analysis of findings. Final write-up and documentation of the project.
%\end{description}

%\begin{table}[htpb]
%	\noindent
%	\begin{tabularx}{\textwidth}{@{}p{3.5cm} X@{}}
%		\textbf{\textsf{Literature Review}}\newline October -- November & Review of relevant literature and assessment of the current state-of-the-art in particle simulations. Additionally, exploration of non-convex global optimization algorithms with corresponding applications and implementations. \\
%		\textbf{\textsf{Algorithmic Design}}\newline December           & Familiarization with AutoPas' tuning strategies and codebase, design of the simulated annealing and genetic algorithm tuning strategies.                                                                                          \\
%		\textbf{\textsf{Pause}}\newline January -- May                  & Pausing work due to an exchange semester during the spring term at UC Berkeley. Work on this project will continue after my return.                                                                \\
%		\textbf{\textsf{Implementation}}\newline June -- July           & Implementation of tuning strategies and integration into the AutoPas framework. Testing and optimization.                                                                                                                                                                                                                                 \\
%		\textbf{\textsf{Documentation}}\newline August                  & Evaluation of the implemented approaches, analysis of findings. Final write-up and documentation of the project.                                                                                                                  \\
%	\end{tabularx}
%\end{table}

\begin{table}[htpb]
	\noindent
	\begin{tabularx}{\textwidth}{@{}p{3.5cm} X@{}}
		\textbf{\textsf{Literature Review}}\newline October -- November & Review of relevant literature and assessment of the current state-of-the-art in particle simulations, with particular emphasis on existing Auto-Tuning approaches. Additionally, exploration of core concepts of non-convex global optimization algorithms, and applications for functions on discrete domains. \\
		\textbf{\textsf{Algorithmic Design}}\newline December & Familiarization with AutoPas' tuning strategies, software architecture and codebase. Design of simulated annealing and genetic algorithm tuning strategies, including definition of necessary parameters, processing of collected samples and integration with the existing tuning framework. \\
		\textbf{\textsf{Pause}}\newline January -- May & Pausing project work due to an exchange semester during the spring term at UC Berkeley. During this period, no active work on the project is planned. Work on the project will resume after returning from the exchange semester. \\
		\textbf{\textsf{Implementation}}\newline June -- July & Implementation of the previously designed tuning strategies and integration into the AutoPas framework. Testing of the implemented algorithms across different simulation scenarios, as well as performance evaluation and optimization of the implementations. \\
		\textbf{\textsf{Documentation}}\newline August & Evaluation of the implemented approaches based on their performance across different simulation scenarios and analysis of strengths and limitations. Final write-up and documentation of implementation and evaluation process including interpretation of results. \\
	\end{tabularx}
\end{table}

\vfill
Niklas Ladurner\hfill \thedate\\[-\baselineskip]

\printbibliography




\end{document}